\section*{Bloque D. Ampliación: Teoría y Aplicaciones Prácticas (70\% / 30\%)}

\noindent\textit{Preguntas 91 a 180. Ampliación del banco oficial con 63 preguntas teóricas y conceptuales (91 a 153) y 27 aplicaciones cuantitativas prácticas (154 a 180).}

\subsection*{Parte I. Preguntas teóricas y conceptuales (91 a 153)}

\q{Cuando una EPS aplica una estructura tarifaria por bloques crecientes (IBT) en el consumo de agua potable, el conjunto presupuestario del hogar:}
\begin{opts}
    \item deja de ser convexo y genera múltiples soluciones de esquina interiores.
    \item sigue siendo convexo y presenta vértices cóncavos hacia el origen.
    \item mantiene una pendiente constante igual a la tarifa media del primer bloque.
    \item se expande hacia afuera en los bloques superiores por efecto del subsidio.
\end{opts}

\q{En presencia de un racionamiento cuantitativo máximo $\bar{x}$ sobre el bien $x$, el conjunto presupuestario se modifica de manera que:}
\begin{opts}
    \item se desplaza paralelamente hacia el origen en toda su extensión lineal.
    \item pierde la propiedad de convexidad al impedir la combinación de cestas.
    \item rota hacia adentro alrededor del intercepto horizontal del bien $y$.
    \item se trunca verticalmente en $x=\bar{x}$ manteniendo la convexidad.
\end{opts}

\q{Si el gobierno otorga un subsidio en especie de $\bar{x}$ unidades de un bien no revendible en lugar de una transferencia monetaria de igual costo fiscal $p_x \bar{x}$, un consumidor racional:}
\begin{opts}
    \item alcanza una utilidad menor o igual que con la transferencia monetaria.
    \item alcanza siempre una utilidad estrictamente mayor por el efecto precio.
    \item alcanza necesariamente la misma utilidad sin importar sus preferencias.
    \item reduce siempre su consumo total del bien subsidiado por debajo de $\bar{x}$.
\end{opts}

\q{Compare un impuesto específico $t$ sobre el consumo del bien $x$ con un impuesto de suma alzada $T$ que recauda exactamente el mismo ingreso fiscal. Para un consumidor con preferencias estrictamente convexas:}
\begin{opts}
    \item el impuesto específico genera una utilidad estrictamente mayor por elasticidad.
    \item ambos impuestos generan exactamente la misma pérdida de utilidad individual.
    \item el impuesto de suma alzada genera una utilidad estrictamente mayor.
    \item el impuesto de suma alzada distorsiona la relación marginal de sustitución.
\end{opts}

\q{Si un consumidor enfrenta un descuento por volumen que reduce el precio marginal del bien $x$ a partir de un umbral $x_0$ (tarifa por bloques decrecientes), entonces el conjunto presupuestario:}
\begin{opts}
    \item deja de ser convexo y puede originar saltos discontinuos en la demanda.
    \item permanece estrictamente convexo garantizando una demanda siempre continua.
    \item se vuelve idéntico al de una tarifa lineal con precio medio ponderado.
    \item impide que el consumidor elija cestas con consumo superior al umbral $x_0$.
\end{opts}

\q{Un impuesto ad valorem a tasa $\tau$ aplicado simultáneamente sobre todos los bienes de la economía es económicamente equivalente a:}
\begin{opts}
    \item un cambio en la pendiente relativa de la recta presupuestaria.
    \item una reducción proporcional del ingreso monetario disponible.
    \item un incremento en la relación marginal de sustitución en el óptimo.
    \item un impuesto progresivo que grava únicamente a los bienes de lujo.
\end{opts}

\q{El teorema de representación de Debreu garantiza la existencia de una función de utilidad continua que representa una relación de preferencias $\succsim$ cuando esta es:}
\begin{opts}
    \item completa, reflexiva y estrictamente convexa sin requerir continuidad.
    \item transitiva, monótona y diferenciable en todo el espacio de bienes.
    \item completa, homotética y aditivamente separable en cada mercancía.
    \item completa, transitiva y continua sobre el conjunto de consumo.
\end{opts}

\q{Las preferencias lexicográficas en $\mathbb{R}^2_+$ no pueden ser representadas por ninguna función de utilidad real porque violan el axioma de:}
\begin{opts}
    \item completitud entre cestas distintas.
    \item transitividad estricta del ordenamiento.
    \item continuidad de las preferencias.
    \item monotonicidad estricta en ambos bienes.
\end{opts}

\q{Si $u(x,y)$ representa las preferencias de un consumidor y $f(\cdot)$ es una transformación estrictamente creciente, entonces la nueva función $v(x,y) = f(u(x,y))$:}
\begin{opts}
    \item multiplica la $RMS$ por la derivada $f'(u)$ en cada punto.
    \item conserva inalterada la $RMS$ en cada cesta del plano.
    \item altera la demanda marshalliana cuando $f$ es no lineal.
    \item preserva necesariamente la concavidad estricta de la utilidad.
\end{opts}

\q{Que una relación de preferencias sea estrictamente convexa implica matemáticamente que cualquier función de utilidad que la represente debe ser:}
\begin{opts}
    \item estrictamente cuasicóncava en todo el conjunto de consumo.
    \item estrictamente cóncava y homogénea de grado uno en bienes.
    \item lineal en al menos uno de los argumentos de la cesta.
    \item estrictamente cuasiconvexa con curvas de indiferencia cóncavas.
\end{opts}

\q{Una característica definitoria de las preferencias homotéticas es que a lo largo de cualquier rayo que parte del origen:}
\begin{opts}
    \item la $RMS$ crece proporcionalmente al nivel de ingreso real.
    \item la elasticidad ingreso de un bien es cero y la del otro supera la unidad.
    \item la $RMS$ permanece constante y las curvas de Engel son lineales.
    \item el gasto en cada bien se mantiene fijo ante variaciones del ingreso.
\end{opts}

\q{Para una función de utilidad cuasilineal $u(x,y) = v(x) + y$ con $v'(x)>0$ y $v''(x)<0$, en soluciones interiores un incremento en el ingreso $m$ provoca que:}
\begin{opts}
    \item el consumo de $x$ aumente en la misma proporción que el ingreso $m$.
    \item el consumo de $y$ permanezca fijo mientras $x$ absorbe el aumento.
    \item ambos bienes se comporten como bienes inferiores de forma simultánea.
    \item el consumo de $x$ no cambie y todo el ingreso extra vaya a $y$.
\end{opts}

\q{En la función de utilidad Stone-Geary $u(x,y) = (x-\gamma_x)^\alpha (y-\gamma_y)^{1-\alpha}$ con $\gamma_x, \gamma_y > 0$, la participación del gasto en cada bien depende del ingreso porque:}
\begin{opts}
    \item las preferencias no son homotéticas debido a los consumos mínimos.
    \item la elasticidad de sustitución entre los bienes es estrictamente cero.
    \item el bien $x$ se convierte en un bien Giffen cuando el ingreso es bajo.
    \item la suma de las participaciones presupuestarias supera la unidad.
\end{opts}

\q{En una función de utilidad CES $u(x,y) = (x^\rho + y^\rho)^{1/\rho}$ con $\rho \le 1$ y $\rho \neq 0$, conforme el parámetro $\rho \to -\infty$ las preferencias convergen a:}
\begin{opts}
    \item sustitutos perfectos lineales con elasticidad $\sigma \to \infty$.
    \item preferencias Cobb-Douglas con elasticidad unitaria $\sigma = 1$.
    \item complementos perfectos de tipo Leontief con $\sigma \to 0$.
    \item preferencias cuasilineales con efecto ingreso nulo en el bien $x$.
\end{opts}

\q{En el problema de maximización de utilidad sujeto a $p_x x + p_y y = m$, el multiplicador de Lagrange $\lambda^*$ evaluado en el óptimo representa:}
\begin{opts}
    \item la relación marginal de sustitución entre los dos bienes.
    \item la utilidad marginal del ingreso monetario en el equilibrio.
    \item el gasto mínimo necesario para alcanzar una utilidad unitaria.
    \item la elasticidad precio cruzada de la demanda marshalliana.
\end{opts}

\q{Con preferencias estrictamente convexas y diferenciables, si en el punto $(0, m/p_y)$ se cumple que $RMS(0, m/p_y) < p_x/p_y$, entonces el óptimo del consumidor:}
\begin{opts}
    \item es necesariamente una solución interior con ambos consumos positivos.
    \item no existe porque la restricción presupuestaria no es acotada.
    \item implica gastar todo el ingreso exclusivamente en el bien $x$.
    \item es una solución de esquina con consumo nulo del bien $x$.
\end{opts}

\q{En un problema de maximización de utilidad con dos bienes y una restricción lineal, la condición suficiente de segundo orden para un máximo interior exige que el determinante del Hessiano orlado $|\bar{H}|$ sea:}
\begin{opts}
    \item estrictamente negativo en el punto crítico.
    \item idénticamente igual a cero en todo el dominio.
    \item estrictamente positivo en el punto crítico.
    \item igual al negativo del multiplicador de Lagrange.
\end{opts}

\q{Si un consumidor tiene preferencias estrictamente cóncavas (no convexas) representadas por $u(x,y) = x^2 + y^2$, su elección óptima ante precios estrictamente positivos:}
\begin{opts}
    \item se ubica siempre en una solución de esquina del presupuesto.
    \item cumple la tangencia interior $RMS = p_x/p_y$ con ambos bienes positivos.
    \item es única e igual a la mitad del ingreso en cada mercancía.
    \item depende únicamente del ingreso y no de los precios relativos.
\end{opts}

\q{Aplicando el teorema de la envolvente al problema primal del consumidor, el efecto de primer orden de un pequeño aumento en $p_x$ sobre la utilidad máxima $V(p,m)$ es:}
\begin{opts}
    \item igual a $+\lambda^* x^*(p,m)$ debido a la reasignación óptima.
    \item nulo en el margen porque el consumidor sustituye hacia el bien $y$.
    \item dependiente exclusivamente de la elasticidad cruzada de demanda.
    \item igual a $-\lambda^* x^*(p,m)$, sin requerir derivar $x^*$ ni $y^*$.
\end{opts}

\q{Cuando las preferencias son localmente no saciadas, la solución óptima del problema primal del consumidor siempre satisface la restricción presupuestaria con igualdad porque:}
\begin{opts}
    \item la función de utilidad debe ser homogénea de grado uno en cantidades.
    \item si sobrara ingreso, existiría una cesta factible estrictamente preferida.
    \item los precios relativos exigen gastar igual proporción en cada bien.
    \item el multiplicador de Lagrange podría tomar valores estrictamente negativos.
\end{opts}

\q{La función de utilidad indirecta $V(p,m)$ satisface la propiedad de homogeneidad de grado cero en $(p,m)$, lo cual significa económicamente que:}
\begin{opts}
    \item no existe ilusión monetaria ante cambios proporcionales de precios e ingreso.
    \item la utilidad máxima se duplica cuando los precios y el ingreso se duplican a la vez.
    \item el consumidor gasta siempre la misma fracción del ingreso en cada bien normal.
    \item las demandas marshallianas son independientes de los precios relativos vigentes.
\end{opts}

\q{Con respecto al vector de precios $p$, la función de utilidad indirecta $V(p,m)$ es siempre:}
\begin{opts}
    \item no decreciente y estrictamente cóncava en $p$.
    \item lineal y homogénea de grado uno en $p$.
    \item estrictamente cuasicóncava en el vector de precios.
    \item no creciente y cuasiconvexa en $p$.
\end{opts}

\q{Si $V(p,m)$ es diferenciable en $(p,m)$ con $\partial V/\partial m > 0$, la demanda marshalliana $x_i(p,m)$ se recupera mediante la identidad de Roy como:}
\begin{opts}
    \item el cociente positivo $\frac{\partial V/\partial p_i}{\partial V/\partial m}$.
    \item el negativo del cociente inverso $\frac{\partial V/\partial m}{\partial V/\partial p_i}$.
    \item el negativo del cociente $\frac{\partial V/\partial p_i}{\partial V/\partial m}$.
    \item la derivada cruzada de segundo orden $\frac{\partial^2 V}{\partial p_i \partial m}$.
\end{opts}

\q{Respecto al vector de precios $p$, la función de gasto $e(p,u)$ es siempre cóncava porque:}
\begin{opts}
    \item la utilidad marginal del ingreso es decreciente en todos los bienes del mercado.
    \item al subir un precio, el consumidor sustituye y gasta menos que con la cesta fija.
    \item el gasto mínimo crece más que proporcionalmente ante el alza de un solo precio relativo.
    \item las preferencias deben ser estrictamente cóncavas para poder minimizar el gasto.
\end{opts}

\q{La función de gasto $e(p,u)$ es homogénea de grado uno en los precios $p$, lo que implica que si todos los precios se multiplican por $t>0$:}
\begin{opts}
    \item el gasto mínimo permanece constante y las demandas hicksianas caen en $1/t$.
    \item las demandas hicksianas se multiplican por $t$ para preservar la utilidad.
    \item la utilidad alcanzable se reduce proporcionalmente en el factor $1/t$.
    \item el gasto mínimo para lograr $u$ se multiplica por $t$ y $h(p,u)$ no varía.
\end{opts}

\q{Según el lema de Shephard en la teoría del consumidor, si $e(p,u)$ es diferenciable respecto a $p_i$, entonces $\frac{\partial e(p,u)}{\partial p_i}$ es igual a:}
\begin{opts}
    \item la demanda hicksiana (compensada) $h_i(p,u)$.
    \item la demanda marshalliana (ordinaria) $x_i(p,m)$.
    \item el negativo de la utilidad marginal del ingreso.
    \item el gasto total asignado al bien $i$ en el óptimo.
\end{opts}

\q{Las identidades fundamentales que vinculan el problema primal y el dual en el óptimo del consumidor establecen que para $u = V(p,m)$ y $m = e(p,u)$:}
\begin{opts}
    \item $x_i(p, e(p,u)) = V(p,m)$ y $h_i(p, V(p,m)) = e(p,u)$.
    \item $x_i(p, e(p,u)) = h_i(p,u)$ y $h_i(p, V(p,m)) = x_i(p,m)$.
    \item $e(p, V(p,m)) = u$ y $V(p, e(p,u)) = m$ para todo $p$.
    \item $\partial x_i/\partial p_i = \partial h_i/\partial p_i$ para todo bien normal.
\end{opts}

\q{La inversa del multiplicador de Lagrange del problema primal ($1/\lambda^*$) coincide exactamente con:}
\begin{opts}
    \item la pendiente de la curva de demanda compensada $\frac{\partial h_i}{\partial p_i}$.
    \item la elasticidad ingreso de la función de utilidad indirecta.
    \item el gasto marginal de la utilidad $\frac{\partial e(p,u)}{\partial u}$.
    \item el costo medio por unidad de utilidad $e(p,u)/u$.
\end{opts}

\q{En la ecuación de Slutsky $\frac{\partial x_i}{\partial p_j} = \frac{\partial h_i}{\partial p_j} - x_j \frac{\partial x_i}{\partial m}$, el término $\frac{\partial h_i}{\partial p_j}$ cuando $i = j$:}
\begin{opts}
    \item es siempre menor o igual que cero por la concavidad de $e(p,u)$.
    \item es positivo únicamente cuando el bien $i$ es un bien Giffen.
    \item es estrictamente negativo solo si el bien $i$ es un bien normal.
    \item depende del signo de la elasticidad ingreso del bien $i$.
\end{opts}

\q{Para que un bien $i$ sea un bien Giffen ($\frac{\partial x_i}{\partial p_i} > 0$), es condición necesaria (aunque no suficiente) que:}
\begin{opts}
    \item sea un bien normal con participación presupuestaria cercana a cero.
    \item tenga un efecto sustitución hicksiano estrictamente positivo en el margen.
    \item su elasticidad precio cruzada respecto a otros bienes sea negativa.
    \item sea inferior y su efecto ingreso supere al efecto sustitución.
\end{opts}

\q{La matriz de sustitución de Slutsky $S(p,m)$ con elementos $s_{ij} = \frac{\partial h_i}{\partial p_j}$ satisface las propiedades de:}
\begin{opts}
    \item asimetría ($s_{ij}=-s_{ji}$) y ser definida positiva.
    \item simetría ($s_{ij}=s_{ji}$) y ser semidefinida negativa.
    \item simetría únicamente cuando todas las preferencias son homotéticas.
    \item tener todos sus elementos fuera de la diagonal mayores que cero.
\end{opts}

\q{Dos bienes $i$ y $j$ se denominan sustitutos netos (o hicksianos) cuando $\frac{\partial h_i}{\partial p_j} > 0$. A diferencia de la sustitución bruta, esta relación es:}
\begin{opts}
    \item posible solo si ambos bienes tienen igual elasticidad ingreso.
    \item incompatible con la ley de la demanda compensada.
    \item necesariamente simétrica entre ambos bienes.
    \item asimétrica cuando los bienes tienen distinta participación del gasto.
\end{opts}

\q{En una economía con solo dos bienes ($x$ e $y$) y preferencias estrictamente convexas, los dos bienes deben ser necesariamente:}
\begin{opts}
    \item complementos netos ($\partial h_x / \partial p_y \le 0$).
    \item sustitutos brutos ($\partial x / \partial p_y > 0$).
    \item bienes normales con elasticidad ingreso unitaria.
    \item sustitutos netos ($\partial h_x / \partial p_y \ge 0$).
\end{opts}

\q{La ley de agregación de Engel establece que la suma de las elasticidades ingreso ponderadas por la participación en el gasto ($\sum_i s_i \varepsilon_{i,m}$) es igual a:}
\begin{opts}
    \item $1$, por lo que no todos los bienes pueden ser inferiores.
    \item $0$, debido a la homogeneidad de grado cero de la demanda ordinaria.
    \item el negativo de la participación del bien numerario en el gasto total.
    \item la elasticidad de sustitución promedio de la cesta de consumo.
\end{opts}

\q{La condición de agregación de Cournot, derivada de diferenciar la restricción presupuestaria respecto al precio $p_j$, establece que:}
\begin{opts}
    \item la suma ponderada $\sum_i s_i \varepsilon_{ij}$ es positiva e igual a $+s_j$.
    \item la suma ponderada $\sum_i s_i \varepsilon_{ij}$ es siempre nula ($0$).
    \item la suma ponderada $\sum_i s_i \varepsilon_{ij}$ equivale a $-s_j$.
    \item la suma simple de elasticidades $\sum_i \varepsilon_{ij}$ es igual a $-1$.
\end{opts}

\q{En forma de elasticidades, la ecuación de Slutsky para el propio precio de un bien $i$ se expresa como $\varepsilon_{ii} = \varepsilon^h_{ii} - s_i \varepsilon_{i,m}$. De esta expresión se deduce que para un bien normal:}
\begin{opts}
    \item la demanda hicksiana es siempre más elástica que la marshalliana.
    \item la demanda marshalliana es más elástica que la hicksiana.
    \item ambas demandas tienen exactamente la misma pendiente si $s_i > 0$.
    \item el signo de $\varepsilon_{ii}$ es positivo cuando $s_i$ es elevado.
\end{opts}

\q{Si un bien es necesario ($0 < \varepsilon_{i,m} < 1$), entonces a medida que aumenta el ingreso del hogar, su curva de Engel (con $m$ en el eje horizontal y $x_i$ en el vertical) es:}
\begin{opts}
    \item creciente y la participación del gasto $s_i$ disminuye con el ingreso.
    \item decreciente y la participación del gasto $s_i$ aumenta con el ingreso.
    \item una recta desde el origen con participación del gasto $s_i$ constante.
    \item creciente y estrictamente convexa con $s_i$ convergiendo a la unidad.
\end{opts}

\q{Si la demanda de agua potable residencial presenta una elasticidad precio propia de $\varepsilon_p = -0.35$ en un tramo dado, un incremento tarifario del 10\% aprobado por SUNASS provocará en el gasto total de los usuarios:}
\begin{opts}
    \item una disminución aproximada del 3.5\% por el efecto sustitución.
    \item una caída aproximada del 6.5\% debido a la pérdida de ingreso real.
    \item ningún cambio en la facturación al compensarse precio y cantidad.
    \item un aumento aproximado del 6.5\% al ser una demanda inelástica.
\end{opts}

\q{En la regulación tarifaria de agua potable en el Perú (SUNASS), la estructura de subsidios cruzados mediante bloques crecientes busca conciliar equidad y eficiencia porque:}
\begin{opts}
    \item elimina completamente la elasticidad ingreso del consumo de agua potable.
    \item protege el consumo básico y acerca al costo marginal los bloques altos.
    \item iguala la tarifa media al precio del primer bloque para todos los usuarios.
    \item maximiza el consumo total de agua en los usuarios industriales y comerciales.
\end{opts}

\q{La función de demanda de mercado de un bien privado se obtiene sumando horizontalmente las demandas individuales. Para que la demanda agregada dependa solo del ingreso total y no de su distribución (forma polar de Gorman):}
\begin{opts}
    \item todos los consumidores deben tener exactamente el mismo nivel de ingreso.
    \item las funciones de utilidad individuales deben ser todas estrictamente cóncavas.
    \item las curvas de Engel individuales deben ser lineales y con igual pendiente.
    \item la elasticidad precio cruzada entre todos los bienes debe ser nula.
\end{opts}

\q{La elasticidad precio de la demanda de agua potable tiende a ser mayor (en valor absoluto) en el largo plazo que en el corto plazo principalmente porque:}
\begin{opts}
    \item el efecto ingreso desaparece completamente después de un año fiscal.
    \item la curva de demanda de largo plazo es necesariamente vertical.
    \item el agua potable se convierte en un bien Giffen al cambiar la tarifa.
    \item los hogares adoptan equipos ahorradores y ajustan sus hábitos.
\end{opts}

\q{La condición de homogeneidad de grado cero de la demanda marshalliana $x_i(p,m)$ implica que la suma de todas las elasticidades precio (propia y cruzadas) más la elasticidad ingreso es igual a:}
\begin{opts}
    \item uno ($\sum_j \varepsilon_{ij} + \varepsilon_{i,m} = 1$).
    \item cero ($\sum_j \varepsilon_{ij} + \varepsilon_{i,m} = 0$).
    \item menos uno ($\sum_j \varepsilon_{ij} + \varepsilon_{i,m} = -1$).
    \item la participación presupuestaria del bien ($s_i$).
\end{opts}

\q{Si un consumidor elige la cesta $x^0$ a precios $p^0$ y la cesta $x^1 \neq x^0$ a precios $p^1$, se produce una violación del Axioma Débil de la Preferencia Revelada (WARP) cuando:}
\begin{opts}
    \item $x^1$ era asequible en $0$ ($p^0 \cdot x^1 \le p^0 \cdot x^0$) y $x^0$ es asequible en $1$.
    \item $x^1$ era inasequible en $0$ ($p^0 \cdot x^1 > p^0 \cdot x^0$) y $x^0$ es inasequible en $1$.
    \item $x^1$ era asequible en $0$ ($p^0 \cdot x^1 \le p^0 \cdot x^0$) pero $x^0$ es inasequible en $1$.
    \item ambas cestas tienen igual costo en $0$ pero $x^0$ cuesta estrictamente más en $1$.
\end{opts}

\q{¿Por qué el cumplimiento de WARP es suficiente para garantizar la existencia de preferencias racionales en un mundo de dos bienes, pero insuficiente con tres o más bienes?}
\begin{opts}
    \item Porque con dos bienes la función de utilidad siempre es homotética.
    \item Porque WARP no requiere que el consumidor gaste todo su ingreso.
    \item Porque con tres o más bienes pueden surgir ciclos intransitivos.
    \item Porque con tres o más bienes la matriz de Slutsky deja de existir.
\end{opts}

\q{El teorema de Afriat establece que un conjunto finito de datos de precios y cantidades $(p^t, x^t)$ es racionalizable por una función de utilidad continua, cóncava y localmente no saciada si y solo si los datos satisfacen:}
\begin{opts}
    \item el Axioma Débil de la Preferencia Revelada (WARP) únicamente.
    \item el Axioma Generalizado de la Preferencia Revelada (GARP).
    \item la condición de homogeneidad lineal de las curvas de Engel.
    \item la constancia de las participaciones presupuestarias en el tiempo.
\end{opts}

\q{Si el índice de cantidades de Laspeyres satisface $L_Q = \frac{p^0 \cdot x^1}{p^0 \cdot x^0} < 1$, el análisis de preferencia revelada permite concluir con certeza que el consumidor:}
\begin{opts}
    \item mejoró inequívocamente su bienestar en el periodo $1$.
    \item es indiferente entre las cestas de ambos periodos fiscales.
    \item violó necesariamente el axioma de monotonicidad estricta.
    \item estaba mejor en el periodo base $0$ que en el periodo $1$.
\end{opts}

\q{Al medir el aumento del costo de vida tras un alza en los precios relativos, el índice de precios de Laspeyres (usado en el IPC) tiende a sobreestimar el verdadero índice de costo de vida de Konüs porque:}
\begin{opts}
    \item fija la cesta base e ignora el efecto sustitución.
    \item evalúa los precios utilizando la cesta final después de sustituir.
    \item excluye por definición a los bienes cuya demanda es inelástica.
    \item asume que todas las preferencias son sustitutos perfectos lineales.
\end{opts}

\q{En contraste con el índice de Laspeyres, el índice de precios de Paasche $P_P = \frac{p^1 \cdot x^1}{p^0 \cdot x^1}$ constituye:}
\begin{opts}
    \item una cota superior del índice verdadero evaluado en la utilidad inicial $u^0$.
    \item una medida exacta del bienestar independientemente de la sustitución.
    \item una cota inferior del índice verdadero evaluado en la utilidad final $u^1$.
    \item un indicador siempre superior al índice de Laspeyres ante cualquier cambio.
\end{opts}

\q{Ante un incremento en el precio de un bien normal de $p_0$ a $p_1$, el ordenamiento correcto en magnitud (valor absoluto) entre la Variación Compensatoria ($|VC|$), la Variación Equivalente ($|VE|$) y el cambio en el Excedente del Consumidor ($|\Delta EC|$) es:}
\begin{opts}
    \item $|VC| < |\Delta EC| < |VE|$ (estrictamente).
    \item $|\Delta EC| < |VE| < |VC|$ (estrictamente).
    \item igualdad exacta entre $|VC|$, $|\Delta EC|$ y $|VE|$.
    \item $|VE| < |\Delta EC| < |VC|$ (estrictamente).
\end{opts}

\q{¿Bajo qué condición sobre las preferencias coinciden exactamente la Variación Compensatoria, la Variación Equivalente y el cambio en el Excedente del Consumidor ante un cambio en el precio del bien $x$?}
\begin{opts}
    \item Cuando la utilidad es cuasilineal en el numerario (efecto ingreso nulo).
    \item Cuando las preferencias son homotéticas de tipo Cobb-Douglas estándar.
    \item Cuando el bien $x$ es un bien inferior con alta participación en el gasto.
    \item Cuando el cambio de precio es compensado por un impuesto ad valorem.
\end{opts}

\q{Ante una reducción del precio del agua potable de $p_0$ a $p_1$, la Variación Compensatoria ($VC$) se define formalmente mediante la función de gasto como:}
\begin{opts}
    \item $VC = e(p_0, u_1) - e(p_1, u_1)$ (monto adicional para alcanzar $u_1$).
    \item $VC = e(p_1, u_1) - e(p_0, u_0)$ (variación del ingreso nominal total).
    \item $VC = e(p_0, u_0) - e(p_1, u_0)$ (monto máximo a pagar por la rebaja).
    \item $VC = V(p_1, m) - V(p_0, m)$ (cambio cardinal en la utilidad indirecta).
\end{opts}

\q{Ante un aumento del precio de un bien inferior (no Giffen), el ordenamiento en valor absoluto entre $|VC|$, $|\Delta EC|$ y $|VE|$ se invierte respecto a un bien normal, cumpliéndose:}
\begin{opts}
    \item $|VE| < |\Delta EC| < |VC|$.
    \item $|VC| < |\Delta EC| < |VE|$.
    \item $|\Delta EC| < |VC| < |VE|$.
    \item $|VC| = |VE| < |\Delta EC|$.
\end{opts}

\q{Al comparar dos políticas tarifarias alternativas $A$ y $B$ desde un mismo equilibrio inicial $(p_0, m)$, la Variación Equivalente ($VE$) es preferida por los economistas sobre la $VC$ porque:}
\begin{opts}
    \item evalúa todas las alternativas a un mismo precio base $p_0$.
    \item no requiere conocer la función de gasto ni las preferencias del hogar.
    \item coincide siempre con la variación del gasto nominal observado en facturas.
    \item elimina el efecto sustitución al mantener constantes las cantidades iniciales.
\end{opts}

\q{El axioma clave de von Neumann-Morgenstern que garantiza que la función de utilidad sobre loterías sea lineal en las probabilidades es el axioma de:}
\begin{opts}
    \item convexidad estricta sobre el espacio de bienes físicos.
    \item homogeneidad de grado uno respecto a la riqueza inicial.
    \item no saciedad local en el conjunto presupuestario determinista.
    \item independencia de alternativas irrelevantes (sustituibilidad).
\end{opts}

\q{A diferencia de la utilidad ordinal bajo certidumbre, una función de utilidad de von Neumann-Morgenstern $u(w)$ es única únicamente hasta:}
\begin{opts}
    \item cualquier transformación estrictamente creciente y diferenciable.
    \item transformaciones homogéneas de grado cero en la riqueza inicial.
    \item transformaciones afines positivas $v(w) = a + b u(w)$ con $b > 0$.
    \item multiplicaciones por una constante negativa con traslación positiva.
\end{opts}

\q{Para un agente estrictamente averso al riesgo ($u''(w) < 0$), frente a una lotería no degenerada $\tilde{w}$ con media $E[\tilde{w}]$, su equivalente cierto $EC$ y su prima de riesgo $\pi$ satisfacen:}
\begin{opts}
    \item $EC > E[\tilde{w}]$ y $\pi = EC - E[\tilde{w}] > 0$.
    \item $EC < E[\tilde{w}]$ y $\pi = E[\tilde{w}] - EC > 0$.
    \item $EC = E[\tilde{w}]$ y $\pi = 0$ por la desigualdad de Jensen.
    \item $EC < E[\tilde{w}]$ y $\pi = u(E[\tilde{w}]) - E[u(\tilde{w})] < 0$.
\end{opts}

\q{El coeficiente de aversión absoluta al riesgo de Arrow-Pratt se define como $r_A(w) = -\frac{u''(w)}{u'(w)}$. Una función de utilidad con aversión absoluta al riesgo constante (CARA) implica que ante un aumento en la riqueza:}
\begin{opts}
    \item la proporción de la riqueza invertida en el activo riesgoso es constante.
    \item el monto invertido en el activo riesgoso crece proporcionalmente a $w$.
    \item el agente se vuelve neutral al riesgo cuando la riqueza supera un umbral.
    \item el monto monetario invertido en un activo riesgoso permanece constante.
\end{opts}

\q{Si un inversionista posee preferencias con aversión relativa al riesgo constante (CRRA), como $u(w) = \frac{w^{1-\gamma}}{1-\gamma}$ con $\gamma > 0$, entonces al duplicarse su riqueza inicial:}
\begin{opts}
    \item mantiene constante la fracción invertida en el activo riesgoso.
    \item mantiene constante el monto absoluto en soles en el activo riesgoso.
    \item reduce a la mitad su disposición a pagar por cualquier seguro justo.
    \item aumenta su coeficiente de aversión absoluta al riesgo $r_A(w)$.
\end{opts}

\q{Un individuo estrictamente averso al riesgo enfrenta una pérdida $L$ con probabilidad $p$ y puede contratar un seguro cuya prima por unidad de cobertura es $q$. Si el seguro es actuarialmente justo ($q = p$), el óptimo del individuo es:}
\begin{opts}
    \item contratar cobertura parcial ($0 < K^* < L$) para ahorrar la prima.
    \item contratar cobertura total ($K^* = L$) igualando su riqueza en ambos estados.
    \item no contratar seguro ($K^* = 0$) porque el valor esperado no cambia.
    \item sobreasegurarse ($K^* > L$) para obtener mayor riqueza si ocurre el siniestro.
\end{opts}

\q{En la ecuación de Slutsky con dotaciones iniciales $(\omega_x, \omega_y)$, un aumento en el precio $p_x$ genera un efecto riqueza adicional que es positivo si el consumidor es:}
\begin{opts}
    \item comprador neto del bien $x$ ($\omega_x < x^*$).
    \item autárquico en ambos bienes ($\omega_x = x^*$).
    \item vendedor neto del bien $x$ ($\omega_x > x^*$).
    \item consumidor de un bien inferior con $\omega_x = 0$.
\end{opts}

\q{En un modelo de elección intertemporal con dos periodos y dotaciones $(m_1, m_2)$, si un consumidor es inicialmente prestamista ($c_1^* < m_1$) y la tasa de interés real $r$ aumenta, entonces por preferencia revelada:}
\begin{opts}
    \item permanece siendo prestamista y su bienestar aumenta estrictamente.
    \item se convierte necesariamente en prestatario por el efecto ingreso.
    \item su consumo presente $c_1^*$ disminuye de forma inequívoca siempre.
    \item su bienestar disminuye si el consumo presente es un bien normal.
\end{opts}

\q{En el modelo de oferta laboral (elección entre ocio y consumo), la curva de oferta de trabajo puede doblarse hacia atrás (pendiente negativa) a salarios elevados cuando:}
\begin{opts}
    \item el ocio es un bien inferior y refuerza al efecto sustitución.
    \item el salario real cae por debajo del salario de reserva del trabajador.
    \item las preferencias entre consumo y ocio son de tipo Cobb-Douglas.
    \item el ocio es normal y el efecto ingreso domina al de sustitución.
\end{opts}

\q{En presencia de imperfecciones en el mercado de crédito donde la tasa de interés para endeudarse ($r_d$) supera a la tasa para ahorrar ($r_a$), la restricción presupuestaria intertemporal presenta:}
\begin{opts}
    \item un conjunto presupuestario no convexo que elimina la autarquía como óptimo.
    \item un quiebre en la dotación, más empinado al endeudarse.
    \item una pendiente única igual al promedio aritmético entre $r_d$ y $r_a$.
    \item una zona de ahorro más empinada que la zona de endeudamiento.
\end{opts}

\subsection*{Parte II. Aplicaciones prácticas y ejercicios cuantitativos (154 a 180)}

\q{Una EPS cobra $S/\ 2$ por $m^3$ para los primeros $10\ m^3$ de agua potable y $S/\ 5$ por cada $m^3$ adicional. Si un hogar cuenta con un presupuesto de $S/\ 70$ para agua ($x$) y otros bienes ($y$, con $p_y = 1$), ¿cuál es la cantidad máxima de agua que puede adquirir si $y = 0$?}
\begin{opts}
    \item $14.0\ m^3$
    \item $25.0\ m^3$
    \item $20.0\ m^3$
    \item $35.0\ m^3$
\end{opts}

\q{Un consumidor con ingreso $m = 120$ enfrenta precios $p_x = 4$ y $p_y = 6$, eligiendo inicialmente la cesta $(x, y) = (15, 10)$. Si se impone un impuesto específico $t = 2$ por unidad de $x$, ¿cuál es la compensación de ingreso de Slutsky necesaria para que la cesta inicial siga siendo exactamente asequible?}
\begin{opts}
    \item $S/\ 30.00$
    \item $S/\ 20.00$
    \item $S/\ 45.00$
    \item $S/\ 60.00$
\end{opts}

\q{Sea $u(x,y) = x^{0.4}y^{0.6}$ con $p_x = 4$, $p_y = 6$ y $m = 200$. La cesta óptima $(x^*, y^*)$ del consumidor es:}
\begin{opts}
    \item $(25, 16)$
    \item $(20, 20)$
    \item $(30, 13)$
    \item $(16, 22)$
\end{opts}

\q{Un hogar tiene preferencias Stone-Geary $u(x,y) = (x - 4)^{0.5}(y - 2)^{0.5}$, donde $x$ representa $m^3$ de agua potable. Si $p_x = 5$, $p_y = 10$ y $m = 100$, la demanda óptima de agua $x^*$ es:}
\begin{opts}
    \item $12.0\ m^3$
    \item $14.0\ m^3$
    \item $16.0\ m^3$
    \item $10.0\ m^3$
\end{opts}

\q{Dado $u(x,y) = 12\sqrt{x} + y$, con precios $p_x = 3$, $p_y = 1$ e ingreso $m = 50$, determine la cesta óptima $(x^*, y^*)$:}
\begin{opts}
    \item $(9, 23)$
    \item $(6, 32)$
    \item $(4, 38)$
    \item $(8, 26)$
\end{opts}

\q{Para la función de utilidad $u(x,y) = \min(3x, 2y)$ con $p_x = 6$, $p_y = 8$ e ingreso $m = 180$, el nivel máximo de utilidad alcanzable $V(p,m)$ es:}
\begin{opts}
    \item $20.0$
    \item $30.0$
    \item $45.0$
    \item $60.0$
\end{opts}

\q{Un consumidor tiene preferencias $u(x,y) = 3x + 5y$ y enfrenta $p_x = 6$, $p_y = 8$ y $m = 120$. La cesta óptima $(x^*, y^*)$ es:}
\begin{opts}
    \item $(0, 15)$
    \item $(20, 0)$
    \item $(10, 15)$
    \item $(12, 6)$
\end{opts}

\q{Dada la función CES $u(x,y) = (x^{0.5} + y^{0.5})^2$, con $p_x = 1$, $p_y = 4$ y $m = 100$, calcule la cantidad óptima $x^*$:}
\begin{opts}
    \item $50.0$
    \item $64.0$
    \item $20.0$
    \item $80.0$
\end{opts}

\q{Si un consumidor maximiza $u(x,y) = x^2 + 2y^2$ sujeto a $2x + 4y = 40$ con $x, y \ge 0$, el valor máximo de la utilidad es:}
\begin{opts}
    \item $200$
    \item $300$
    \item $400$
    \item $600$
\end{opts}

\q{La función de utilidad indirecta de un consumidor es $V(p_x, p_y, m) = \frac{m}{2p_x + 3p_y}$. Si $p_x = 2$, $p_y = 2$ y $m = 50$, aplique la identidad de Roy para hallar $x^*$:}
\begin{opts}
    \item $10.0$
    \item $12.5$
    \item $15.0$
    \item $20.0$
\end{opts}

\q{Si la función de utilidad indirecta es $V(p_x, p_y, m) = \ln(m) - 0.3\ln(p_x) - 0.7\ln(p_y)$, ¿cuánto vale la demanda marshalliana $x^*$ cuando $p_x = 9$ y $m = 300$?}
\begin{opts}
    \item $21.0$
    \item $10.0$
    \item $30.0$
    \item $15.0$
\end{opts}

\q{Para un consumidor con utilidad $u(x,y) = x^{0.5}y^{0.5}$, ingreso $m = 120$ y precios $p_x = 4$, $p_y = 9$, el valor de la función de utilidad indirecta $V(p_x, p_y, m)$ es:}
\begin{opts}
    \item $12.0$
    \item $15.0$
    \item $20.0$
    \item $10.0$
\end{opts}

\q{Dada la función de utilidad $u(x,y) = x^{0.5}y^{0.5}$, calcule el gasto mínimo $e(p_x, p_y, u)$ necesario para alcanzar $u = 5$ cuando $p_x = 4$ y $p_y = 16$:}
\begin{opts}
    \item $S/\ 80.00$
    \item $S/\ 64.00$
    \item $S/\ 96.00$
    \item $S/\ 40.00$
\end{opts}

\q{Sea la función de gasto $e(p_x, p_y, u) = 3u p_x^{1/3} p_y^{2/3}$. Usando el lema de Shephard, halle la demanda hicksiana $h_x(p_x, p_y, u)$ evaluada en $p_x = 1$, $p_y = 8$ y $u = 10$:}
\begin{opts}
    \item $20.0$
    \item $30.0$
    \item $40.0$
    \item $60.0$
\end{opts}

\q{Un consumidor tiene función de gasto $e(p_x, p_y, u) = u(\sqrt{p_x} + \sqrt{p_y})^2$. Evaluada en $p_x = 4$, $p_y = 9$ y $u = 4$, la demanda compensada $h_x$ es:}
\begin{opts}
    \item $12.0$
    \item $10.0$
    \item $15.0$
    \item $18.0$
\end{opts}

\q{Para $u(x,y) = x^{0.5}y^{0.5}$ con $m = 160$ y $p_y = 4$, el precio de $x$ sube de $p_x^0 = 4$ a $p_x^1 = 16$. Según la descomposición de Slutsky, los efectos sustitución ($ES$) e ingreso ($EI$) sobre $x$ son respectivamente:}
\begin{opts}
    \item $ES = -10.0$ y $EI = -5.0$
    \item $ES = -5.0$ y $EI = -10.0$
    \item $ES = -12.5$ y $EI = -2.5$
    \item $ES = -7.5$ y $EI = -7.5$
\end{opts}

\q{En un consumidor con $u(x,y) = x^{0.5}y^{0.5}$, $m = 160$ y $p_y = 4$, si $p_x$ aumenta de $4$ a $16$, ¿cuánto vale el efecto sustitución de Hicks ($ES_H$) sobre el bien $x$?}
\begin{opts}
    \item $-7.5$ unidades
    \item $-12.5$ unidades
    \item $-10.0$ unidades
    \item $-15.0$ unidades
\end{opts}

\q{Un bien representa el $25\%$ del presupuesto familiar ($s_x = 0.25$), posee una elasticidad precio compensada $\varepsilon^h_{xx} = -0.40$ y una elasticidad ingreso $\varepsilon_{x,m} = 1.20$. Su elasticidad precio marshalliana $\varepsilon_{xx}$ es:}
\begin{opts}
    \item $-0.10$
    \item $-0.70$
    \item $-0.55$
    \item $-0.85$
\end{opts}

\q{En una canasta de dos bienes, el agua potable ($x$) absorbe el $20\%$ del gasto y tiene una elasticidad ingreso de $0.40$. Por la ley de agregación de Engel, la elasticidad ingreso del resto de bienes ($y$) es:}
\begin{opts}
    \item $1.15$
    \item $0.60$
    \item $1.25$
    \item $1.40$
\end{opts}

\q{Para una canasta de dos bienes con participaciones $s_x = 0.40$ y $s_y = 0.60$, si la elasticidad precio propia del bien $x$ es $\varepsilon_{xx} = -0.70$, entonces por la condición de Cournot la elasticidad cruzada $\varepsilon_{yx}$ vale:}
\begin{opts}
    \item $-0.30$
    \item $-0.40$
    \item $-0.12$
    \item $-0.20$
\end{opts}

\q{La demanda residencial de agua potable de una localidad sigue la ecuación $q = 60 - 2p$ (en $m^3$ mensuales). Cuando la tarifa es $p = 10$, la elasticidad precio de la demanda es:}
\begin{opts}
    \item $-0.33$
    \item $-1.00$
    \item $-0.50$
    \item $-2.00$
\end{opts}

\q{Un consumidor con preferencias cuasilineales $u(x,y) = 20\sqrt{x} + y$ ($p_y = 1$, $m = 100$) enfrenta un alza en $p_x$ de $2$ a $5$. La pérdida de bienestar medida por $|VC| = |VE| = |\Delta EC|$ es:}
\begin{opts}
    \item $S/\ 30.00$
    \item $S/\ 25.00$
    \item $S/\ 40.00$
    \item $S/\ 50.00$
\end{opts}

\q{Para $u(x,y) = x^{0.5}y^{0.5}$ con $m = 100$ y $p_y = 1$, si $p_x$ sube de $1$ a $4$, ¿cuáles son los valores absolutos de la Variación Compensatoria ($|VC|$) y la Variación Equivalente ($|VE|$)?}
\begin{opts}
    \item $|VC| = 50$ y $|VE| = 100$
    \item $|VC| = 100$ y $|VE| = 50$
    \item $|VC| = 75$ y $|VE| = 75$
    \item $|VC| = 150$ y $|VE| = 50$
\end{opts}

\q{Un agente con utilidad $u(w) = \sqrt{w}$ enfrenta una lotería que paga $S/\ 100$ o $S/\ 400$ con igual probabilidad ($0.5$). Su equivalente cierto ($EC$) y prima de riesgo ($\pi$) son:}
\begin{opts}
    \item $EC = 250$ y $\pi = 0$
    \item $EC = 200$ y $\pi = 50$
    \item $EC = 150$ y $\pi = 100$
    \item $EC = 225$ y $\pi = 25$
\end{opts}

\q{Para un individuo con función de utilidad logarítmica $u(w) = \ln(w)$ y riqueza $w = 50$, los coeficientes de aversión absoluta ($r_A$) y relativa ($r_R$) al riesgo de Arrow-Pratt son:}
\begin{opts}
    \item $r_A = 0.02$ y $r_R = 1.00$
    \item $r_A = 1.00$ y $r_R = 0.02$
    \item $r_A = 0.04$ y $r_R = 2.00$
    \item $r_A = 0.50$ y $r_R = 1.00$
\end{opts}

\q{Un agente con utilidad $u(x,y) = x^{0.5}y^{0.5}$ posee una dotación inicial $(\omega_x, \omega_y) = (12, 12)$ y enfrenta precios $p_x = 3$, $p_y = 1$. En el óptimo, el agente:}
\begin{opts}
    \item compra $4$ unidades de $x$ y vende $12$ de $y$.
    \item vende $4$ unidades de $x$ y compra $12$ de $y$.
    \item vende $6$ unidades de $x$ y compra $18$ de $y$.
    \item permanece en autarquía sin intercambiar bienes.
\end{opts}

\q{Un consumidor con preferencias $u(c_1, c_2) = c_1^{0.5}c_2^{0.5}$ recibe ingresos $m_1 = 120$ y $m_2 = 88$, y enfrenta una tasa de interés real $r = 10\%$. En el equilibrio intertemporal:}
\begin{opts}
    \item se endeuda en $20$ en el periodo $1$ y consume $c_2^* = 66$.
    \item ahorra $10$ en el periodo $1$ y consume $c_2^* = 99$.
    \item ahorra $20$ en el periodo $1$ y consume $c_2^* = 110$.
    \item consume exactamente su ingreso en ambos periodos ($c_1^* = 120$).
\end{opts}
